% GENERATED by md2tex.py from ../draft_v5_en/back/H_appendix_training.md -- do not edit by hand.
% Change the Markdown and run `python md2tex.py`; edits here are overwritten.

\chapter{Training specification: detail}
\label{app:H}

\section{Reward function}
\label{sec:H.1}

The reward is the same on both tracks. It is calculated on the \textbf{ground-truth} state plus progress, so it does not depend on the observation:

\begingroup
\small
\begin{verbatim}
r = w_fwd · max(progress, 0)
  - w_ey  · |ey|
  - w_eps · |epsi|
  - w_ds  · |Δsteering|
  - w_term · [terminated_off_road]
\end{verbatim}
\endgroup

where \texttt{progress} is the normalised progress along the centre line, \texttt{Δsteering} is the change in the policy's raw steering (not the value after the cage; \S\ref{sec:7.2.2}), and \texttt{[terminated\_\allowbreak{}off\_\allowbreak{}road]} is 1 only if the episode ends because the vehicle left the road (not because of a \cagerule{05} emergency; \S\ref{sec:7.2.4}). The nominal weights (which may be tuned experimentally) are:

% layout: table, natural, \small, 131 of 165 mm
\begin{table}[!htbp]
\caption{Nominal reward weights.}\label{tab:H.1}
\small\setlength{\tabcolsep}{4pt}
\begin{tabular}{@{}lll@{}}
  \toprule
  \tabhead{Parameter} & \tabhead{Value} & \tabhead{Rationale} \\
  \midrule
  \texttt{w\_fwd} (forward\_progress) & 1.0 & Rewards real progress ($\approx$1.0/step at cruise) \\
  \texttt{w\_ey} (lateral\_error) & 2.5 & Main penalty: lateral offset \\
  \texttt{w\_eps} (heading\_error) & 0.75 & Secondary penalty: heading \\
  \texttt{w\_ds} (steer\_delta) & 0.20 & Actuation smoothness (over the raw $\Delta$steering, v1.2; \S\ref{sec:7.2.2}) \\
  \texttt{w\_term} (termination) & 25.0 & Discourages leaving the road \\
  \bottomrule
\end{tabular}
\end{table}

The forward term uses normalised progress, not speed. Since speed is fixed, a \texttt{w\_fwd·speed} term would be a constant that does not tell behaviours apart, and it left \texttt{explained\_\allowbreak{}variance ≈ 0}. The high termination penalty (25.0) gives priority to staying on the \textbf{road}. It only applies when the vehicle leaves the road: a \cagerule{05} emergency ends the episode with no penalty, because the cage's intervention is part of the dynamics and not a punishment (\S\ref{sec:7.2.4}). The weights are \texttt{[provisional, M-P1..M-P4]}. Details are in \texttt{docs/\allowbreak{}10\_\allowbreak{}reward\_\allowbreak{}function.\allowbreak{}md}.

\section{Hyperparameters}
\label{sec:H.2}

The table shows the full effective configuration. The E-main column is the camera run \texttt{ppo\_\allowbreak{}newcam\_\allowbreak{}complex\_\allowbreak{}b\_\allowbreak{}2024\_\allowbreak{}1M}, and the F baseline is the state run \texttt{ppo\_\allowbreak{}train\_\allowbreak{}2024\_\allowbreak{}200k}.

% layout: table, natural, \footnotesize, 159 of 165 mm
\begin{table}[!htbp]
\caption{Effective hyperparameters of the E-main and F-baseline runs.}\label{tab:H.2}
\footnotesize\setlength{\tabcolsep}{4pt}
\begin{tabular}{@{}llll@{}}
  \toprule
  \tabhead{Parameter} & \tabhead{E-main (camera)} & \tabhead{Baseline (F, state)} & \tabhead{Source / note} \\
  \midrule
  \texttt{policy} & CnnPolicy & MlpPolicy & policy network \\
  \texttt{total\_timesteps} & 1,000,000 (planned; stopped at $\approx$662k) & 200,000 & \texttt{[provisional, M-P7]} \\
  \texttt{learning\_rate} & 3$\times$10\textsuperscript{$-$4}, linear anneal & 3$\times$10\textsuperscript{$-$4} constant & E: \texttt{lr\_schedule: linear} \\
  \texttt{target\_kl} & 0.5 & — (no brake) & E: trust region brake (\S\ref{sec:7.3}) \\
  \texttt{normalize\_reward} & True (\texttt{VecNormalize}) & False & E: stabilises the critic (\S\ref{sec:7.3}) \\
  \texttt{clip\_range\_vf} & 0.2 & null & E: value clip over the normalised reward \\
  \texttt{gamma} & 0.99 & 0.99 & = SB3 default \\
  \texttt{n\_steps} & 1,024 & 1,024 & $\approx$ 1 episode \\
  \texttt{batch\_size} & 64 & 64 & = SB3 default \\
  \texttt{n\_epochs} & 10 & 10 & SB3 default \\
  \texttt{gae\_lambda} & 0.95 & 0.95 & SB3 default \\
  \texttt{clip\_range} & 0.2 & 0.2 & SB3 default \\
  \texttt{ent\_coef} & 0.0 & 0.0 & no entropy bonus \\
  \texttt{vf\_coef} & 0.5 & 0.5 & SB3 default \\
  \texttt{max\_grad\_norm} & 0.5 & 0.5 & SB3 default \\
  \texttt{device} & auto (CUDA if present) & cpu & E: the CNN benefits from a GPU \\
  \bottomrule
\end{tabular}
\end{table}

The four stability settings of the E-main run (\texttt{target\_kl}, linear LR decay, \texttt{VecNormalize(\allowbreak{}norm\_\allowbreak{}reward)} and \texttt{clip\_range\_vf}) are not in the F baseline. They were added after seeing that PPO with a CNN and visual randomization is much less stable than with the state vector (\S\ref{sec:7.3}). \texttt{norm\_obs} is left False, so evaluation and inference are not affected and \texttt{ep\_rew\_mean} in the curve stays unnormalised (comparable with the baseline). The camera budget is at least 1M steps, and a pilot of about 20k steps checks that the loop works before that budget is spent.

\section{Comparative study of algorithms and checkpoints}
\label{sec:H.3}

The table follows the chain run $\rightarrow$ checkpoint $\rightarrow$ evaluation of the later study. All evaluation values come from the nominal scenario with randomization turned off. The peak shown is the peak of the training curve and does not always match the saved checkpoints, so the final choice is made by closed-loop evaluation.

% layout: table, fixed, \footnotesize, 164 of 165 mm
\begin{table}[!htbp]
\caption[Chain run $\rightarrow$ checkpoint $\rightarrow$ evaluation of the later SAC study.]{Chain run $\rightarrow$ checkpoint $\rightarrow$ evaluation of the later SAC study. The intervention percentages in monitoring are counterfactual activations: they are logged, but the cage's action is not applied.}\label{tab:H.3}
\footnotesize\setlength{\tabcolsep}{4pt}
\begin{tabular}{@{}>{\RaggedRight\arraybackslash}p{39.6mm}>{\RaggedRight\arraybackslash}p{38.0mm}>{\RaggedRight\arraybackslash}p{18.2mm}>{\RaggedRight\arraybackslash}p{24.7mm}>{\RaggedRight\arraybackslash}p{31.7mm}@{}}
  \toprule
  \tabhead{Action $\cdot$ configuration} & \tabhead{Training evidence} & \tabhead{Checkpoint evaluated} & \tabhead{\scn{NOM-01}, enforcement} & \tabhead{\scn{NOM-01}, monitoring} \\
  \midrule
  \textbf{1-D}, SAC \texttt{auto}, seed 2024 $\cdot$ \texttt{sac\_\allowbreak{}newcam\_\allowbreak{}complex\_\allowbreak{}b\_\allowbreak{}2024\_\allowbreak{}1M} & peak 720.0 @ 89,089; manual stop at 307,201 of the 1M planned & 75k (\texttt{58631022…}) & 5.12 laps; 19.8 mm; 0 emerg.; 48.3\,\% \cagerule{06} & 5.13 laps; 23.3 mm; 0 emerg. \\
  \addlinespace[2pt]
  \textbf{1-D}, \texttt{ent\_coef=0.005}, seed 2024 $\cdot$ \texttt{sac\_\allowbreak{}newcam\_\allowbreak{}entfix\_\allowbreak{}complex\_\allowbreak{}b\_\allowbreak{}2024\_\allowbreak{}1M} & peak 722.5 @ 82,945; without the abrupt collapse; stop at 260,097 & 75k (\texttt{b74505ac…}) & 5.04 laps; 21.6 mm; 0 emerg.; 9.1\,\% \cagerule{06} & 5.04 laps; 21.6 mm; 0 emerg. \\
  \addlinespace[2pt]
  \textbf{1-D}, \texttt{ent\_coef=0.005}, seed 42 $\cdot$ \texttt{sac\_\allowbreak{}newcam\_\allowbreak{}entfix\_\allowbreak{}complex\_\allowbreak{}b\_\allowbreak{}42\_\allowbreak{}120k} & peak 744.3 @ 87,041; replica bounded to 120,833 & 75k (\texttt{4d09e43c…}) & 4.63 laps; \textbf{12.3 mm}; 0 emerg.; \textbf{2.3\,\% \cagerule{06}} & \textbf{pending: no nominal \texttt{\_mon} run exists} \\
  \addlinespace[2pt]
  \textbf{1-D}, \texttt{ent\_coef=0.005}, seed 666 $\cdot$ \texttt{sac\_\allowbreak{}newcam\_\allowbreak{}entfix\_\allowbreak{}complex\_\allowbreak{}b\_\allowbreak{}666\_\allowbreak{}120k} & peak 606.9 @ 80,897; replica bounded to 120,833 & 75k (\texttt{18c80fce…}) & 5.00 laps; 14.0 mm; 0 emerg.; 5.3\,\% \cagerule{06} & 5.00 laps; 14.0 mm; 0 emerg.; 6.2\,\% \cagerule{06} counterfactual \\
  \addlinespace[2pt]
  \textbf{1-D}, \texttt{ent\_coef=0.005}, buffer 200k, seed 2024 $\cdot$ \texttt{sac\_\allowbreak{}newcam\_\allowbreak{}entfix\_\allowbreak{}buf200\_\allowbreak{}2024\_\allowbreak{}180k} & sustained band 690–745; peak 744.7 @ 155,649; stop at 180,225 & 150k (\texttt{a5c5f3c4…}) & 4.94 laps; 26.9 mm; 0 emerg.; 14.4\,\% \cagerule{06} & not executed \\
  \addlinespace[2pt]
  \textbf{2-D}, SAC \texttt{auto}, seed 2024 $\cdot$ \texttt{sac\_\allowbreak{}gz2d\_\allowbreak{}tuned\_\allowbreak{}complex\_\allowbreak{}b\_\allowbreak{}2024\_\allowbreak{}1M} & collapse–\allowbreak{}recovery cycles; peak 527.0 @ 153,601; stop at 250,881 & edge 175k (\texttt{e8934d51…}) & 3.45 laps; 34.8 mm; 1 \cagerule{05} stop & 4.31 laps; 32.3 mm; 0 emerg. \\
  \addlinespace[2pt]
  \textbf{2-D}, \texttt{ent\_coef=0.005}, seed 2024 $\cdot$ \texttt{sac\_\allowbreak{}gz2d\_\allowbreak{}tuned\_\allowbreak{}entfix\_\allowbreak{}2024\_\allowbreak{}1M} & peak 558.7 @ 77,825; rise without abrupt cycles; stop at 176,129 & 75k (\texttt{b76724c7…}) & \textbf{4.32 laps; 17.1 mm; 0 emerg.}; 17.1\,\% \cagerule{06} & 4.31 laps; 16.3 mm; 0 emerg. \\
  \addlinespace[2pt]
  \textbf{2-D}, \texttt{ent\_coef=0.005}, seed 42 $\cdot$ \texttt{sac\_\allowbreak{}gz2d\_\allowbreak{}tuned\_\allowbreak{}entfix\_\allowbreak{}42\_\allowbreak{}120k} & peak 270.9 @ 47,105; replica bounded to 120,833 & 50k (\texttt{cbde3836…}) & \textbf{4.97 laps; 18.2 mm; 0 emerg.}; 46.4\,\% \cagerule{06} & 4.84 laps; 22.6 mm; 39 steps with a counterfactual \cagerule{05} trigger \\
  \bottomrule
\end{tabular}
\end{table}
